\section{Extend the stored path into complete Huffman inflation}
Use the attachment helper in \path{drafts/png-huffman-construction/attach.py} on your
current stored-stage program, with a fresh output path. It replaces the two supplied
Huffman development labels, adds read-only format tables and appends internal helper
scaffolds. Your completed stored, checksum, collector and loader regions remain intact.
An altered handler boundary needs manual reconciliation; the helper refuses to guess
which instructions it may replace. Preparing a named recipe changes controls rather
than completing a student region.

The format constants follow \href{https://www.rfc-editor.org/rfc/rfc1951.html}{RFC 1951}.
The tables supplied with this exercise are immutable data: the code-length permutation
and the bases and extra-bit counts for lengths and distances. Per-block counts, ordered
symbols and decoded length vectors belong to the existing call frame. Rebuilding a
block must not inherit a previous block's alphabet. The supplied calls establish table
arguments; the following edits develop the algorithms behind those calls.

\Needspace{12\baselineskip}
\begin{longtable}{p{40mm}p{23mm}p{85mm}}
  \toprule Private region & Frame offset & Capacity and representation                             \\\midrule
  Literal/length table    & 160          & 16 native word counts, then 288 native halfword symbols \\
  Distance table          & 800          & 16 word counts, then 32 halfword symbols                \\
  Code-length table       & 928          & 16 word counts, then 19 halfword symbols                \\
  Placement cursors       & 1032         & 16 native words, reused by each table build             \\
  Combined length vector  & 1096         & 320 bytes for the two data alphabets                    \\
  Code-length lengths     & 1416         & 19 bytes, indexed by code-length symbol                 \\\bottomrule
\end{longtable}
The capacities fit the 2048-byte frame already supplied. All offsets are relative to
\texttt{\$fp}; they are illustration details rather than part of the completed entry
contract. A table begins with sixteen counts indexed by code length zero through fifteen.
Its ordered symbols start 64 bytes later. Length zero means absent. The placement
cursors are scratch: sorting one alphabet changes them, and the next build reuses them.

Build the fixed lengths first, then the shared table builder and symbol reader. Add
literals and overlapping matches before reading dynamic descriptions. An unfinished
helper uses the owning inflater's development exit, preserving its caller's saved
state. Do not enter an unfinished helper without that frame. Fixed-literal streams
can start succeeding before matches and dynamic headers are implemented; that is a
checkpoint of construction, not evidence of a completed inflater.
